% 通用工作稿，不冒充目标期刊官方模板。由 render.py 实际套用。
% !TEX program = xelatex
\documentclass[12pt,a4paper]{ctexart}
\usepackage[a4paper,margin=2.54cm]{geometry}
\usepackage{amsmath,amssymb,graphicx,booktabs,tabularx,array}
\usepackage{caption,float}
\usepackage[hidelinks]{hyperref}
\setCJKmainfont{SimSun}
\setCJKsansfont{SimHei}
\setmainfont{Times New Roman}
\linespread{1.5}
\captionsetup{font=small,labelformat=empty}
\setlength{\parindent}{2em}
\title{\sffamily\bfseries\fontsize{22}{28}\selectfont 合单发车对即时配送效率与时效的影响：事件驱动仿真教学算例}
\author{}
\date{}
\begin{document}
\maketitle
教学演示声明：本文为 Idea2Paper 技能演示（synthetic teaching fixture），数据为合成数据，结论不构成实证研究，不用于真实投稿。

\begin{abstract}
即时配送中，订单在取货点的合单等待窗口（batching window）同时影响配送成本与订单时效，二者存在权衡。本文以事件驱动仿真教学算例考察该权衡：在单一取货点、欧氏距离、固定送达期限（到达后 25 分钟）与固定随机种子下，比较逐单发车（S1）与合单发车（S2，窗口 w=5 分钟）两种策略，覆盖到达率 λ∈\{0.2, 0.5, 1.0\}（单/分钟）。结果表明：合单显著降低配送成本——车辆批次分别减少 68.0\%、78.1\%、85.1\%，总行驶距离分别减少 22.7\%、27.5\%、30.2\%；代价是平均订单等待升至 3.1–7.6 分钟，准时率由约 98–99\% 降至 0–4.9\%。在 λ=0.5 下扫描窗口 w∈\{2,5,8\}，准时率随 w 单调下降（10.5\%→0.95\%→0.32\%），距离持续微降，确认"成本—时效"权衡。本文仅为教学算例：合成数据、单取货点、无限车辆、贪心路径等简化假设下，不声称实证发现或算法创新。
\textbf{关键词：} 即时配送；合单发车；事件驱动仿真；教学算例
\end{abstract}
\section{引言}
即时配送（外卖、同城快递）中，平台在订单到达与派车之间做合单决策：等待更多订单再发车可摊薄每单行驶里程与车辆占用，但增加订单等待并威胁时效。该问题属于带时间窗的取送货问题（PDPTW）的运营层变体，Ropke 与 Pisinger 提出自适应大邻域搜索（ALNS）求解静态 PDPTW [1]。近年研究进一步将预测与决策结合（predict-then-optimize）以利用订单数据改进决策，如 Bertsimas 与 Kallus 的 prescriptive analytics 框架 [2] 与 Elmachtoub 与 Grigas 的 SPO+ 方法 [3]。

本算例不做方法创新，而是用完全可复现的教学算例展示"合单窗口"这一机制的核心权衡，作为技能演示与后续真实研究的脚手架。边界：合成数据、单取货点、无限车辆、贪心路径，结论限于设定范围内。

\section{问题与模型}
\textbf{设定。} 单一取货点位于原点；订单按泊松过程到达（速率 λ 单/分钟），目的地均匀分布于 [−5,5]×[−5,5]（km）；仿真时长 T=600 分钟；订单送达期限为到达时刻后 25 分钟；配送速度 0.6 km/分钟，每单服务时间 3 分钟。

\textbf{策略。} S1 逐单发车：订单到达即刻派车，一车一单。S2 合单发车：订单在取货点排队，自队列首个订单到达起计时，满窗口 w 分钟（或积满 B=10 单）发车，一车承运当前队列全部订单，路径按与原点距离升序贪心串联。

\textbf{指标。} 车辆批次数 N\_batch（车辆占用代理）、总行驶距离 D\_total（km，含返回原点）、平均订单等待 wait\_mean（分钟）、准时率 ontime（送达时刻 ≤ 期限的订单占比）。

\textbf{正确性检查。} w=0 时 S2 应完全退化为 S1（同一订单序列下 N\_batch 与 D\_total 一致，容差 1e-9）。

\section{方法与算法}
事件驱动仿真：先用固定种子生成各 λ 的订单序列（写入 data/demand.json，登记为合成数据），同一序列供两种策略公平比较；随后对每个 (λ, 策略) 组合执行确定性仿真（无随机噪声）。运行由 skill 的 runner 记录输入输出哈希与退出码，结果存 results/logistics.json 与 results/logistics.csv。实现仅用 Python 标准库，种子由运行环境注入（RESEARCH\_SEED=42）。

\section{实验与结果}
\begin{table}[htbp]
\centering
\caption*{表 1：λ 扫描（S1 逐单 vs S2 合单 w=5，同一订单序列，种子 42）}
\begin{tabular}{cccccccc}
\toprule
λ（单/分钟） & 订单数 & 批次 S1→S2 & 批次节省 & 距离 S1→S2 (km) & 距离节省 & 等待 S2 (分) & 准时率 S1→S2 \\
\midrule
0.2 & 122 & 122→39 & 68.0\% & 905.9→700.3 & 22.7\% & 7.58 & 99.2\%→4.9\% \\
0.5 & 315 & 315→69 & 78.1\% & 2360.5→1711.5 & 27.5\% & 4.00 & 98.1\%→0.95\% \\
1.0 & 596 & 596→89 & 85.1\% & 4600.5→3210.8 & 30.2\% & 3.11 & 98.5\%→0\% \\
\bottomrule
\end{tabular}
\end{table}
图 1 与图 2 分别给出车辆批次与平均等待随到达率的变化：合单在到达密度越高时节省越大（批次节省 68\%→85\%），而平均等待随密度上升反而下降（7.6→3.1 分钟，因窗口内积单更快）。

\begin{figure}[htbp]
\centering
\includegraphics[width=\linewidth,height=0.8\textheight,keepaspectratio]{\detokenize{figures/fig2.png}}
\caption*{图 1：车辆批次随订单到达率的变化（S1 vs S2, w=5）}
\end{figure}
注：合成教学数据；种子42，仿真时长600分钟；S1为逐单发车，S2为合单发车；批次单位为次。

\begin{figure}[htbp]
\centering
\includegraphics[width=\linewidth,height=0.8\textheight,keepaspectratio]{\detokenize{figures/fig3.png}}
\caption*{图 2：平均订单等待随订单到达率的变化（S1 vs S2, w=5）}
\end{figure}
注：合成教学数据；与图1使用相同订单序列；等待单位为分钟。

图 3 在 λ=0.5 下扫描窗口 w∈\{2,5,8\}：窗口越大距离越低（1788.4→1711.5→1614.0 km），但准时率单调下降（10.5\%→0.95\%→0.32\%），平均等待升至 5.82 分钟——"成本—时效"权衡清晰可见。

\begin{figure}[htbp]
\centering
\includegraphics[width=\linewidth,height=0.8\textheight,keepaspectratio]{\detokenize{figures/fig4.png}}
\caption*{图 3：准时率随合单窗口的变化（λ=0.5, S2）}
\end{figure}
注：合成教学数据；种子42；准时率为满足期限的订单量/订单总量（\%），窗口单位为分钟。

\section{讨论与结论}
\textbf{权衡解释。} 合单通过把多个订单合并为一次取货点出发的行程摊薄里程与车辆占用；到达密度越高、窗口越长，合并机会越多，成本节省越大。但等待直接消耗送达期限：本算例期限仅 25 分钟，w≥2 即显著恶化准时率；S2 在 w=5、λ=0.2 时准时率仅 4.9\%，说明"合单保成本"在紧时效设定下代价极高。

\textbf{正确性与可信度。} w=0 退化检查通过（S2 与 S1 数值一致）；claims 绑定结果 JSON 字段并在容差 1e-6 内；三篇文献 DOI 经 Crossref 元数据核验（未声称全文阅读）。

\textbf{局限（教学性质）。} 合成数据；单取货点与无限车辆忽略容量与骑手可用性；贪心路径非最优；未做敏感性（除 w）与统计推断；不声称实证发现或新算法。真实研究需接入真实订单数据、考虑车辆容量与动态指派、并用解析或最优基线对照。

\textbf{结论。} 本算例确认即时配送合单存在"成本—时效"权衡：合单可降低 22.7–30.2\% 行驶距离与 68–85\% 车辆批次，但代价为等待与准时率，权衡由窗口与期限共同决定。扩展方向包括容量约束下的动态合单与 predict-then-optimize 式窗口决策 [2,3]。

\section{参考文献}
[1] Ropke, S., \& Pisinger, D. (2006). An adaptive large neighborhood search heuristic for the pickup and delivery problem with time windows. Transportation Science, 40(4), 455-472.

[2] Bertsimas, D., \& Kallus, N. (2020). From predictive to prescriptive analytics. Management Science, 66(3), 1025-1044.

[3] Elmachtoub, A. N., \& Grigas, P. (2022). Smart predict-then-optimize. Operations Research, 70(5), 2644-2662.

\section{附录：数据与代码可得性}
合成数据（data/demand.json）、仿真代码（code/logistics.py）、结果（results/）与图（figures/）随项目一并提供，登记为合成教学数据（urn:synthetic:logistics-demo-v1），许可为教学用途；本演示不涉及真实数据、伦理审批或作者利益冲突。AI 工具使用说明见 availability.md。

\end{document}
